Commutation of the real derivatives gives Hermitian symmetry. In a diagonal entry the terms for $\alpha\ne\beta$ cancel in pairs, and each term with $\alpha=\beta$ has coefficient one. This gives the trace. Applying the Hessian map to $\xi\xi^{\mathsf t}$ gives $gg^*$. Every nonnegative real form is a sum of such forms. The ordinary product rule gives the identity for $H(uv)$. In four slots, the real transpose of a quaternionic unitary $U$ acts on the encoded covector by $U^*$. The rank-one calculation yields $H(u\circ U)=U^*(H(u)\circ U)U$ and linearity extends the calculation to any real Hessian. Real orthogonal matrices on the coordinate index preserve the three-slot subspace, giving the covariance there.

Write $P_a=p_a+\mathbf k q_a$, so $2\partial_{\bar P_a}=\partial_{p_a}+\mathbf k\partial_{q_a}$. Since $\mathbf{ki}=\mathbf j$, the gradient is $2(\partial_{\bar P_a}u)\mathbf i$. In the Hessian the four coefficients for derivatives $(p_ap_b,p_aq_b,q_ap_b,q_aq_b)$ are $1,-\mathbf k,\mathbf k,1$, giving $4\partial_{\bar P_a}\partial_{P_b}u$. For holomorphic $F$, the complex chain rule gives
\[
\partial_{\bar P_a}\partial_{P_b}\sum_\nu|F_\nu|^2
 =\sum_\nu\overline{J_{\nu a}}\,J_{\nu b},\qquad
G(\rho\circ F)=2J^*(\partial_{\bar F_\nu}\rho)_\nu\,\mathbf i,
\]
where the derivatives of $\rho$ are evaluated at $F(P)$.
